\documentclass[11pt]{article}
\usepackage[T1]{fontenc}
\usepackage{textcomp}
\usepackage[margin=0.75in]{geometry}
\usepackage{amsmath,amssymb,amsthm}
\usepackage{array,booktabs}
\usepackage{hyperref}
\setlength{\parindent}{0pt}
\newtheorem{theorem}{Theorem}
\theoremstyle{definition}
\newtheorem{definition}{Definition}
\title{Flow Proof Artifact: Int}
\author{Generated by \texttt{flow doc proof}}
\date{Flow Proof Book}
\begin{document}
\maketitle
Integer arithmetic definitions extending the natural numbers.\\[0.5em]
\textbf{Source.} landau — *Foundations of Analysis*, Ch. 1\\[0.5em]
\bigskip
\noindent{\large \textbf{Definition 1.} \textit{Adding zero on the left leaves an integer unchanged} \hfill the integers \cdot addition \cdot zero is the left identity}\par
\smallskip\noindent\textit{Coordinate: the integers \cdot addition \cdot zero is the left identity}\par
\smallskip\noindent\textit{Source: landau}\par
\medskip
\noindent\textbf{Goal.} Adding zero on the left leaves an integer unchanged\par
\noindent$\displaystyle \forall m \in \mathbb{Z}\quad 0 + m = m$\par
\medskip
\noindent
\begin{tabular}{@{} >{\raggedright\arraybackslash}p{0.06\textwidth} >{\raggedright\arraybackslash}p{0.47\textwidth} >{\raggedleft\arraybackslash}p{0.41\textwidth} @{}}
\textbf{\#} & \textbf{Proof} & \textbf{Mathematics} \\
\midrule
\textbf{1.} & We stipulate zero is the left identity for addition on the integers — this is a definition, not a derived fact. & \\[0.45em]
\textbf{2.} & This follows directly from the definition: 0 plus m equals m. Hence proven. & $\displaystyle 0 + m = m$ \\[0.45em]
\bottomrule
\end{tabular}
\label{thm:Int-addition-zero_is_the_left_identity}
\par\medskip\hrule
\bigskip
\noindent{\large \textbf{Definition 2.} \textit{Adding zero on the right leaves an integer unchanged} \hfill the integers \cdot addition \cdot zero is the right identity}\par
\smallskip\noindent\textit{Coordinate: the integers \cdot addition \cdot zero is the right identity}\par
\smallskip\noindent\textit{Source: landau}\par
\medskip
\noindent\textbf{Goal.} Adding zero on the right leaves an integer unchanged\par
\noindent$\displaystyle \forall m \in \mathbb{Z}\quad m + 0 = m$\par
\medskip
\noindent
\begin{tabular}{@{} >{\raggedright\arraybackslash}p{0.06\textwidth} >{\raggedright\arraybackslash}p{0.47\textwidth} >{\raggedleft\arraybackslash}p{0.41\textwidth} @{}}
\textbf{\#} & \textbf{Proof} & \textbf{Mathematics} \\
\midrule
\textbf{1.} & We stipulate zero is the right identity for addition on the integers — this is a definition, not a derived fact. & \\[0.45em]
\textbf{2.} & This follows directly from the definition: m plus 0 equals m. Hence proven. & $\displaystyle m + 0 = m$ \\[0.45em]
\bottomrule
\end{tabular}
\label{thm:Int-addition-zero_is_the_right_identity}
\par\medskip\hrule
\bigskip
\noindent{\large \textbf{Derived fact 3.} \textit{Negating twice returns the original integer} \hfill the integers \cdot negation \cdot double negation returns the value}\par
\smallskip\noindent\textit{Coordinate: the integers \cdot negation \cdot double negation returns the value}\par
\smallskip\noindent\textit{Source: landau}\par
\medskip
\noindent\textbf{Goal.} Negating twice returns the original integer\par
\noindent$\displaystyle \forall m \in \mathbb{Z}\quad -(-m) = m$\par
\medskip
\noindent
\begin{tabular}{@{} >{\raggedright\arraybackslash}p{0.06\textwidth} >{\raggedright\arraybackslash}p{0.47\textwidth} >{\raggedleft\arraybackslash}p{0.41\textwidth} @{}}
\textbf{\#} & \textbf{Proof} & \textbf{Mathematics} \\
\midrule
\textbf{1.} & We prove that double negation returns the value for negation on the integers. & \\[0.45em]
\textbf{2.} & We can deduce that -(-m) equals m. Hence proven. & $\displaystyle -(-m) = m$ \\[0.45em]
\bottomrule
\end{tabular}
\label{thm:Int-negation-double_negation_returns_the_value}
\par\medskip\hrule
\end{document}
